\chapter{Polynomial time, reductions, and NP}
\label{chap:polytime}

This chapter covers deterministic and nondeterministic time classes, the
polynomial-time function class $\cc{FP}$, polynomial-time many-one (Karp)
reductions, and the theory of $\cc{NP}$-completeness built on satisfiability.
All class definitions are formalized, together with the closure of $\cc{P}$ and
$\cc{FP}$ under the usual operations, Cobham's machine-independent
characterization of $\cc{FP}$ in both directions, the direction ``polynomial-time
verifier implies $\cc{NP}$'' of the witness characterization, and the
Cook--Levin theorem with its corollaries for $3\mathrm{SAT}$ and for the
complement of $\mathrm{SAT}$. The main open directions are the converse witness
direction (every $\cc{NP}$ language has a polynomial-time verifier, hence
$\cc{NP} = \Sigma^p_1$), a catalog of further $\cc{NP}$-complete problems with
explicit codecs, and $\cc{FP}$ combinators that would let the existing bespoke
machines be re-proved as short programs through Cobham's theorem.

\section{Time classes}

\begin{definition}[$\cc{DTIME}$]
  \label{def:dtime}
  \lean{Complexity.DTIME}
  \leanok
  \uses{def:language, def:tm-decides}
  For $T \colon \mathbb{N} \to \mathbb{N}$, $\cc{DTIME}(T)$ is the set of
  languages $L$ for which there are a deterministic machine with some number
  $k$ of work tapes and a function $f$ with $f = O(T)$ such that the machine
  decides $L$ within time $f$: on every input $x$ it halts within $f(|x|)$
  steps with $1$ in the first output cell if $x \in L$ and $0$ otherwise. Here
  $f = O(T)$ means that $f(n) \le C \cdot T(n)$ for some constant $C$ and all
  sufficiently large $n$.
\end{definition}

\begin{definition}[$\cc{NTIME}$]
  \label{def:ntime}
  \lean{Complexity.NTIME}
  \leanok
  \uses{def:language, def:ntm-decides}
  $\cc{NTIME}(T)$ is the set of languages $L$ for which some nondeterministic
  machine (two transition functions selected by a choice bit) and some
  $f = O(T)$ satisfy: every computation path on input $x$ halts within
  $f(|x|)$ steps, and $x \in L$ if and only if some choice sequence of length
  $f(|x|)$ ends in the halting state with output $1$.
\end{definition}

\begin{definition}[$\cc{P}$]
  \label{def:P}
  \lean{Complexity.P}
  \leanok
  \uses{def:dtime}
  $\cc{P} = \bigcup_{k} \cc{DTIME}(n^k)$.
\end{definition}

\begin{definition}[$\cc{NP}$]
  \label{def:NP}
  \lean{Complexity.NP}
  \leanok
  \uses{def:ntime}
  $\cc{NP} = \bigcup_{k} \cc{NTIME}(n^k)$. This is the machine definition;
  the certificate form is Theorem~\ref{thm:polytime-fnp-np} (one direction,
  formalized) and Theorem~\ref{thm:polytime-np-verifier} (the other, planned).
\end{definition}

\begin{definition}[$\cc{coNP}$]
  \label{def:coNP}
  \lean{Complexity.coNP}
  \leanok
  \uses{def:NP}
  $\cc{coNP} = \{L : \overline{L} \in \cc{NP}\}$, where $\overline{L}$ is the
  complement of $L$ in the set of all bit strings. The library forms every
  complement class with the same constructor
  (\texttt{Complexity.complClass}).
\end{definition}

\begin{definition}[$\cc{FP}$]
  \label{def:FP}
  \lean{Complexity.FP}
  \leanok
  \uses{def:tm-computes}
  $\cc{FP}$ is the set of functions $f \colon \{0,1\}^* \to \{0,1\}^*$ for which
  some deterministic machine computes $f$ within time $T$ with
  $T = O(n^d)$ for some $d$: on every input $x$ it halts within $T(|x|)$ steps
  with $f(x)$ written on the output tape.
\end{definition}

\begin{definition}[$\cc{EXP}$ and $\cc{NEXP}$]
  \label{def:EXP}
  \lean{Complexity.EXP, Complexity.NEXP}
  \leanok
  \uses{def:dtime, def:ntime}
  $\cc{EXP} = \bigcup_{k} \cc{DTIME}(2^{n^k})$ and
  $\cc{NEXP} = \bigcup_{k} \cc{NTIME}(2^{n^k})$.
\end{definition}

\begin{theorem}[Deterministic inside nondeterministic time]
  \label{thm:polytime-P-subset-NP}
  \lean{Complexity.DTIME_subset_NTIME, Complexity.P_subset_NP, Complexity.P_subset_EXP, Complexity.NP_subset_NEXP, Complexity.EXP_subset_NEXP}
  \leanok
  \uses{def:dtime, def:ntime, def:P, def:NP, def:EXP}
  For every $T$, $\cc{DTIME}(T) \subseteq \cc{NTIME}(T)$. Consequently
  $\cc{P} \subseteq \cc{NP}$, $\cc{P} \subseteq \cc{EXP}$,
  $\cc{NP} \subseteq \cc{NEXP}$, and $\cc{EXP} \subseteq \cc{NEXP}$.
\end{theorem}
\begin{proof}
  \leanok
  A deterministic machine is a nondeterministic machine whose two transition
  functions coincide, and it decides the same language in the same time. The
  exponential inclusions follow from $n^k \le 2^{n^k}$ and monotonicity of
  $\cc{DTIME}$ and $\cc{NTIME}$ under $O(\cdot)$.
\end{proof}

\begin{theorem}[$\cc{P}$, $\cc{NP}$, and $\cc{coNP}$]
  \label{thm:polytime-coNP-structure}
  \lean{Complexity.P_subset_NP_inter_coNP, Complexity.NP_eq_coNP_of_P_eq_NP, Complexity.P_ne_NP_of_NP_ne_coNP}
  \leanok
  \uses{def:P, def:NP, def:coNP}
  $\cc{P} \subseteq \cc{NP} \cap \cc{coNP}$. If $\cc{P} = \cc{NP}$ then
  $\cc{NP} = \cc{coNP}$; equivalently, $\cc{NP} \neq \cc{coNP}$ implies
  $\cc{P} \neq \cc{NP}$.
\end{theorem}
\begin{proof}
  \leanok
  \uses{thm:polytime-P-subset-NP, prop:polytime-P-boolean}
  $\cc{P}$ is closed under complement, so $\cc{P} \subseteq \cc{coNP}$. Under
  $\cc{P} = \cc{NP}$ the class $\cc{NP}$ inherits closure under complement, and
  a class closed under complement equals its complement class.
\end{proof}

\begin{theorem}[$\cc{NP} \subseteq \cc{PSPACE}$]
  \label{thm:polytime-np-subset-pspace}
  \uses{def:NP, def:PSPACE, def:EXP}
  $\cc{NP} \subseteq \cc{PSPACE}$, and hence $\cc{NP} \subseteq \cc{EXP}$.
\end{theorem}
\begin{proof}
  \uses{thm:polytime-np-sigma1, thm:randomness-ph-subset-pspace}
  Two routes are available. Directly: prove
  $\cc{NTIME}(T) \subseteq \cc{NSPACE}(T)$ exactly as
  $\cc{DTIME}(T) \subseteq \cc{DSPACE}(T)$ is proved (a head moves at most one
  cell per step), then apply $\cc{NPSPACE} \subseteq \cc{PSPACE}$
  (\texttt{Complexity.NPSPACE\_subset\_PSPACE}, formalized in the space
  chapter). Alternatively, $\cc{NP} \subseteq \Sigma^p_1 \subseteq \cc{PH}
  \subseteq \cc{PSPACE}$. In both cases $\cc{PSPACE} \subseteq \cc{EXP}$
  (\texttt{Complexity.PSPACE\_subset\_EXP}) gives the second claim.
\end{proof}

\begin{theorem}[Regular languages are in $\cc{P}$]
  \label{thm:polytime-regular-in-P}
  \lean{Complexity.mem_DTIME_of_isRegular, Complexity.mem_P_of_isRegular, Complexity.mem_P_of_nfa,
    Complexity.mem_P_of_twoWayNA, Complexity.mem_P_of_finite_monoid}
  \leanok
  \uses{def:dtime, def:P}
  Every regular binary language, in the sense shared by Mathlib and CSLib, is in
  $\cc{DTIME}(n + 2)$ and hence in $\cc{P}$. In particular, the languages of
  CSLib's finite nondeterministic automata, one-way or two-way, and the
  preimages of subsets of finite monoids under monoid homomorphisms from binary
  strings are in $\cc{P}$.
\end{theorem}
\begin{proof}
  \leanok
  A finite-state scanner machine runs the deterministic automaton over the
  input in one pass and writes whether the final state accepts. For the other
  models, CSLib's subset construction, its two-way-to-one-way theorem, and its
  finite-monoid characterization first show that the language is regular. CSLib's closure properties of regular languages therefore
  give $\cc{P}$-membership of the resulting languages at no further cost.
\end{proof}

\begin{theorem}[Time classes on CSLib's multi-tape machines]
  \label{thm:polytime-cslib-multitape}
  \lean{Complexity.TM.DecidesInTime.decidableInTimeAndSpace,
    Complexity.decidableInTimeAndSpace_of_mem_DTIME, Complexity.decidableInTimeAndSpace_of_mem_P,
    Complexity.TM.DecidesInTimeSpace.decidableInTimeAndSpace,
    Complexity.decidableInTimeAndSpace_of_mem_DTISP,
    Complexity.TM.ComputesInTime.computableInTimeAndSpace,
    Complexity.computableInTimeAndSpace_of_mem_FP}
  \leanok
  \uses{def:dtime, def:P, def:FP}
  If a machine with $k$ work tapes decides $L$ within time $f$, then a
  multi-tape machine in CSLib's model decides $L$ within time and space
  $(2k + 3)(2f + 5)$. CSLib's model has a binary alphabet, a read-only input
  tape clamped to the input, two-way work tapes that start blank, and emitted
  output, and it counts the work cells visited as space. Consequently every
  language in $\cc{DTIME}(T)$ is decided there within time and space $O(T)$,
  and every language in $\cc{P}$ within polynomial time and space. The
  simulation is also space-faithful: a machine deciding $L$ within time $T$
  and space $S$ yields a CSLib decider within time $2T + 4$ and space
  $(2k + 3)(S + 2)$, so $\cc{DTISP}(T, S)$ transfers with time $O(T)$ and space
  $O(S + 1)$. Functions computed within time $T$ are computed there within time
  $3T + 4$, so every $\cc{FP}$ function is CSLib-computable in polynomial time.
\end{theorem}
\begin{proof}
  \leanok
  Each of the original machine's work tapes and its output tape is simulated
  by a data tape and a marker tape whose heads move together; the simulator's
  first step writes a $1$ at position $0$ of every marker tape, standing in
  for $\triangleright$. It then executes one step of the original machine per
  step of its own. A further work tape counts how far the input head has run
  past the input, and a flag in the state distinguishes the blank before the
  input from the blank after it. When the original machine halts, the
  simulator rewinds its copy of the output tape, reads the verdict in cell
  $1$, and emits it as one bit. Rewinding costs at most the original running
  time, and each work head visits at most one cell per step. Under a space
  bound, every simulated head stays within the original machine's head
  bounds. To compute a function, the simulator instead copies output cells
  $1, 2, \ldots$ to CSLib's output until the first blank; only visited cells
  can be nonblank, so the output has length at most $T$.
\end{proof}

\begin{theorem}[CSLib's polynomial time is $\cc{P}$]
  \label{thm:polytime-cslib-P-eq}
  \lean{Complexity.decidesInTime_of_decidableInTimeAndSpace,
    Complexity.mem_DTIME_of_decidableInTimeAndSpace,
    Complexity.mem_P_iff_decidableInTimeAndSpace}
  \leanok
  \uses{thm:polytime-cslib-multitape}
  A CSLib multi-tape machine deciding $L$ within time $t$ yields one of our
  machines deciding $L$ within time $3t$. Consequently a language is in
  $\cc{P}$ if and only if some multi-tape machine in CSLib's model decides it
  within polynomial time.
\end{theorem}
\begin{proof}
  \leanok
  \uses{thm:polytime-cslib-multitape}
  Theorem~\ref{thm:polytime-cslib-multitape} gives one direction. For the
  converse, CSLib's machines are binary, so their tape symbols (blank, $0$,
  $1$) are already ours. Fold each two-way work tape onto a one-sided tape,
  sending CSLib cell $z \ge 0$ to cell $2z + 1$ and $z < 0$ to cell $-2z$, and
  detect a crossing of the fold by bumping into $\triangleright$. Track the
  clamped input head directly and write emitted bits from output cell $1$.
  Each CSLib step takes three of our steps.
\end{proof}

\section{Closure properties of P and FP}

The projections of the pairing codec, and the pairing of two $\cc{FP}$
functions, are in $\cc{FP}$; these facts are stated in the chapter on
encodings. Pairing a value with the unchanged input is
\texttt{Complexity.mem\_FP\_pairWithInput}.

\begin{proposition}[$\cc{P}$ is a Boolean algebra]
  \label{prop:polytime-P-boolean}
  \lean{Complexity.P_compl, Complexity.P_union, Complexity.P_inter, Complexity.P_diff, Complexity.complClass_P}
  \leanok
  \uses{def:P}
  $\cc{P}$ is closed under complement, union, intersection, and set difference.
  In particular $\{L : \overline{L} \in \cc{P}\} = \cc{P}$.
\end{proposition}
\begin{proof}
  \leanok
  Complement runs the decider and flips the output bit, in $2f(n) + 4$ steps.
  Union runs both deciders in sequence, giving
  $\cc{DTIME}(T_1) \cup \cc{DTIME}(T_2) \subseteq \cc{DTIME}(T_1 + T_2)$ in the
  sense of \texttt{Complexity.DTIME\_union}. Intersection and difference follow
  by De Morgan.
\end{proof}

\begin{lemma}[Identity and composition in $\cc{FP}$]
  \label{lem:polytime-FP-comp}
  \lean{Complexity.id_mem_FP, Complexity.mem_FP_comp}
  \leanok
  \uses{def:FP}
  The identity function is in $\cc{FP}$, and if $f, g \in \cc{FP}$ then
  $g \circ f \in \cc{FP}$.
\end{lemma}
\begin{proof}
  \leanok
  The identity is computed by copying the input to the output in $n + 2$
  steps. For composition, the two machines run as a pipeline on disjoint tapes,
  the output of the first becoming the input of the second; since
  $|f(x)| \le T_f(|x|)$, the total time is polynomial.
\end{proof}

\begin{theorem}[Closure under polynomial-time preimages]
  \label{thm:polytime-preimage}
  \lean{Complexity.mem_P_preimage, Complexity.mem_NP_preimage}
  \leanok
  \uses{def:P, def:NP, def:FP}
  If $f \in \cc{FP}$ and $L \in \cc{P}$, then
  $f^{-1}(L) = \{x : f(x) \in L\} \in \cc{P}$. The same holds with $\cc{NP}$ in
  place of $\cc{P}$.
\end{theorem}
\begin{proof}
  \leanok
  Run the transducer for $f$ and then the decider for $L$ on its output, in a
  pipeline of disjoint tapes, deterministically or nondeterministically. The
  output of $f$ has polynomial length, so the second stage runs in polynomial
  time in $|x|$.
\end{proof}

\begin{lemma}[Deciding by a verdict function]
  \label{lem:polytime-decision-fn}
  \lean{Complexity.mem_P_of_decisionFn, Complexity.mem_P_of_decisionFn_bool}
  \leanok
  \uses{def:P, def:FP}
  If $f \in \cc{FP}$ and, for all $x$, $x \in L$ exactly when $f(x)$ contains
  a $1$, then $L \in \cc{P}$. In particular, if $g$ is Boolean-valued,
  $x \mapsto [g(x)]$ is in $\cc{FP}$, and $x \in L \iff g(x) = 1$, then
  $L \in \cc{P}$.
\end{lemma}
\begin{proof}
  \leanok
  \uses{thm:polytime-preimage}
  $L$ is the preimage under $f$ of the polynomial-time language of strings
  containing a $1$. Together with Cobham's theorem this lets membership in
  $\cc{P}$ be established by writing a function rather than a machine.
\end{proof}

\section{Cobham's theorem: FP as a programming language}

\begin{definition}[Cobham's function algebra]
  \label{def:polytime-cobham}
  \lean{Complexity.Cobham, Complexity.CobhamFP, Complexity.smash, Complexity.recNotation}
  \leanok
  Bit strings are read least-significant bit first. Cobham's algebra is the
  smallest class of functions $(\{0,1\}^*)^n \to \{0,1\}^*$, over all arities
  $n$, that contains the projections, the empty-string constant, the bit
  successors $x \mapsto b\,x$ (prepend $b$), and
  $\mathrm{smash}(x, y) = 1^{|x| \cdot |y|}$, and is closed under composition
  and limited recursion on notation: if $g, h_0, h_1, j$ are in the class and
  the function $f$ defined by $f(\varepsilon, v) = g(v)$ and
  $f(b\,x, v) = h_b(x, f(x, v), v)$ satisfies $|f(x, v)| \le |j(x, v)|$ for
  all $x, v$, then $f$ is in the class. The
  unary fragment is $\cc{CobhamFP}$.
\end{definition}

\begin{theorem}[Cobham]
  \label{thm:cobham}
  \lean{Complexity.CobhamFP_eq_FP, Complexity.CobhamFP_subset_FP, Complexity.FP_subset_CobhamFP}
  \leanok
  \uses{def:polytime-cobham, def:FP}
  $\cc{CobhamFP} = \cc{FP}$.
\end{theorem}
\begin{proof}
  \leanok
  \uses{lem:polytime-iterate, lem:polytime-FP-comp}
  Soundness is an induction over the constructors of the algebra; limited
  recursion on notation is a fold whose width clamp is vacuous under Cobham's
  side condition, run by Lemma~\ref{lem:polytime-iterate}. Completeness
  simulates a polynomial-time machine inside the algebra: a configuration is
  one block-aligned string with each tape split at its head, the transition
  function is a finite table, the run is iterated under a clock built from
  $\mathrm{smash}$, and the output is read after a rewind.
\end{proof}

\begin{theorem}[Cobham's theorem at every arity]
  \label{thm:polytime-cobham-arity}
  \lean{Complexity.Cobham.cobham_iff_FPn}
  \leanok
  \uses{def:polytime-cobham, def:FP}
  For every $n$ and $f \colon (\{0,1\}^*)^n \to \{0,1\}^*$, $f$ is in Cobham's
  algebra if and only if there is $g \in \cc{FP}$ with
  $g(\mathrm{enc}(v)) = f(v)$ for all $v$, where $\mathrm{enc}$ is the
  canonical fixed-arity tuple encoding (\texttt{Cobham.encodeVec}).
\end{theorem}
\begin{proof}
  \leanok
  \uses{thm:cobham}
  The soundness induction is carried out at every arity; completeness
  composes the unary simulation with the tuple encoding, which is itself in
  the algebra.
\end{proof}

\begin{lemma}[Iteration in $\cc{FP}$]
  \label{lem:polytime-iterate}
  \lean{Complexity.Cobham.iterate_mem_FP}
  \leanok
  \uses{def:FP}
  Let $F, \mathit{init}, \mathit{ruler}, \mathit{width} \in \cc{FP}$, and
  suppose $|F^m(\mathit{init}(z))| \le |\mathit{width}(z)|$ for every $z$ and
  every $m \le |\mathit{ruler}(z)|$. Then
  $z \mapsto F^{|\mathit{ruler}(z)|}(\mathit{init}(z))$ is in $\cc{FP}$.
\end{lemma}
\begin{proof}
  \leanok
  One machine iterates a truncated step once per bit of its own input, with a
  unary counter that stops applying $F$ after $|\mathit{ruler}(z)|$ rounds;
  the width bound keeps every intermediate state polynomially long. The
  declaration lives in the proof-internal
  module \texttt{Classes/P/Cobham/Internal.lean}; it is the main lever for
  writing polynomial-time loops and is a candidate for promotion to a surface
  module.
\end{proof}

\begin{lemma}[Range concatenation and streaming folds in $\cc{FP}$]
  \label{lem:polytime-fp-combinators}
  \uses{def:FP}
  The following combinators preserve $\cc{FP}$.
  \begin{enumerate}
    \item Bounded concatenation over a range: if
      $(z, i) \mapsto f(z, i)$ is in $\cc{FP}$ (on paired inputs, with $i$ in
      unary) and $m \in \cc{FP}$, then
      $z \mapsto f(z, 0) \, f(z, 1) \cdots f(z, |m(z)| - 1)$ is in $\cc{FP}$.
    \item Streaming finite-state fold: for a finite state set $Q$, a transition
      $\delta \colon Q \times \{0,1\} \to Q$, an initial state, and an output
      map $o \colon Q \times \{0,1\} \to \{0,1\}^*$, the transducer that scans
      $z$ from left to right and emits $o(q_i, z_i)$ at each position is in
      $\cc{FP}$; more generally the state may be a polynomially bounded string
      updated by an $\cc{FP}$ step.
  \end{enumerate}
\end{lemma}
\begin{proof}
  \uses{thm:cobham, lem:polytime-iterate}
  Both are instances of Lemma~\ref{lem:polytime-iterate}, with the state
  carrying the output produced so far, or can be written directly by limited
  recursion on notation. These two combinators describe the Cook--Levin
  emitter (clause families are concatenations over index ranges), the Tseitin
  transducer (one streaming pass over a CNF encoding), and the $\mathrm{SAT}$
  verifier (a streaming evaluation); see the open directions at the end of this
  chapter.
\end{proof}

\section{Witnesses, FNP, and TFNP}

\begin{definition}[Pair language and balanced relations]
  \label{def:pairlang}
  \lean{Complexity.pairLang, Complexity.PolyBalanced}
  \leanok
  \uses{def:pair, def:language}
  For a relation $R \subseteq \{0,1\}^* \times \{0,1\}^*$,
  $\mathrm{pairLang}(R) = \{\langle x, y \rangle : R(x, y)\}$ using the
  library's pairing codec. $R$ is polynomially balanced if there is a
  polynomial $p$ with $R(x, y) \Rightarrow |y| \le p(|x|)$.
\end{definition}

\begin{definition}[$\cc{FNP}$ and $\cc{TFNP}$]
  \label{def:FNP}
  \lean{Complexity.FNP, Complexity.TFNP}
  \leanok
  \uses{def:pairlang, def:P}
  $\cc{FNP}$ is the set of polynomially balanced relations $R$ with
  $\mathrm{pairLang}(R) \in \cc{P}$. $\cc{TFNP}$ is the set of $R \in \cc{FNP}$
  such that every $x$ has some $y$ with $R(x, y)$.
\end{definition}

\begin{lemma}[Guess and verify]
  \label{lem:polytime-guess-verify}
  \lean{Complexity.mem_NP_of_poly_witness}
  \leanok
  \uses{def:NP, def:P, def:pair}
  Let $p$ be a polynomial and $L_0 \in \cc{P}$ with
  $\langle x, y \rangle \in L_0 \Rightarrow |y| \le p(|x|)$. If, for all $x$,
  $x \in L \iff \exists y\, \langle x, y \rangle \in L_0$, then
  $L \in \cc{NP}$.
\end{lemma}
\begin{proof}
  \leanok
  \uses{thm:polytime-preimage}
  The guess-and-verify machine built for $\mathrm{SAT}$ is generic: it guesses
  a string of length at most $|x| + 1$, pairs it with the input, and runs the
  deterministic verifier, which handles linearly bounded witnesses. A general
  polynomial bound is reduced to the linear case by padding the input with a
  polynomial-length ruler and pulling back along the padding map. This lemma
  currently lives in \texttt{Classes/PCP/Internal/GuessVerifyGeneric.lean}.
\end{proof}

\begin{theorem}[$\cc{FNP}$ witnesses give $\cc{NP}$]
  \label{thm:polytime-fnp-np}
  \lean{Complexity.NP.mem_NP_of_FNP, Complexity.NP.witnessLang_mem_NP, Complexity.NP.witnessNTMConstruction}
  \leanok
  \uses{def:FNP, def:NP}
  If $R \in \cc{FNP}$ and $x \in L \iff \exists y\, R(x, y)$ for all $x$, then
  $L \in \cc{NP}$. In particular $\{x : \exists y\, R(x, y)\} \in \cc{NP}$ for
  every $R \in \cc{FNP}$.
\end{theorem}
\begin{proof}
  \leanok
  \uses{lem:polytime-guess-verify}
  The decider of $\mathrm{pairLang}(R)$ is the verifier and the balance
  polynomial bounds the witnesses, so Lemma~\ref{lem:polytime-guess-verify}
  applies. The construction is packaged as the named interface
  \texttt{NP.WitnessNTMConstruction} and discharged unconditionally.
\end{proof}

\begin{theorem}[$\cc{NP}$ languages have polynomial-time verifiers]
  \label{thm:polytime-np-verifier}
  \uses{def:NP, def:FNP}
  For every $L \in \cc{NP}$ there is $R \in \cc{FNP}$ with
  $x \in L \iff \exists y\, R(x, y)$ for all $x$. Together with
  Theorem~\ref{thm:polytime-fnp-np},
  $\cc{NP} = \{L : \exists R \in \cc{FNP}\ \forall x\ (x \in L \iff \exists y\,
  R(x, y))\}$.
\end{theorem}
\begin{proof}
  \uses{thm:cobham, lem:polytime-decision-fn}
  Take as witness the choice sequence of an accepting path, of length
  $f(|x|)$. The verifier runs the machine deterministically along the given
  choices; \texttt{NTM.choiceTM} already does this, and the same run is
  simulated inside Cobham's algebra for the Sipser--Lautemann theorem, so the
  verdict can be written as a Cobham function and converted to $\cc{P}$ with
  Lemma~\ref{lem:polytime-decision-fn}.
\end{proof}

\begin{theorem}[$\cc{NP} = \Sigma^p_1$]
  \label{thm:polytime-np-sigma1}
  \uses{def:NP, def:coNP, def:PH}
  $\cc{NP} = \Sigma^p_1$ and $\cc{coNP} = \Pi^p_1$, for the
  certificate-quantifier levels of Definition~\ref{def:PH}.
\end{theorem}
\begin{proof}
  \uses{lem:polytime-guess-verify, thm:polytime-np-verifier, prop:randomness-ph-levels, prop:polytime-P-boolean}
  $\Sigma^p_1$ is the bounded existential closure of $\cc{P}$ (formalized as
  \texttt{Complexity.SigmaP\_one}). For $\Sigma^p_1 \subseteq \cc{NP}$, apply
  Lemma~\ref{lem:polytime-guess-verify} to $L' \cap \{\langle x, w \rangle :
  |w| \le p(|x|)\}$, whose length test is polynomial-time. For
  $\cc{NP} \subseteq \Sigma^p_1$, use Theorem~\ref{thm:polytime-np-verifier}.
  The $\Pi^p_1$ statement follows by complementation.
\end{proof}

\begin{theorem}[$\cc{NP} \cap \cc{coNP}$ search is total]
  \label{thm:polytime-tfnp}
  \lean{Complexity.orRelation_mem_TFNP_of_NP_coNP_witnesses}
  \leanok
  \uses{def:FNP}
  Let $R_1, R_2 \in \cc{FNP}$ and $L$ satisfy
  $x \in L \iff \exists y\, R_1(x, y)$ and
  $x \notin L \iff \exists y\, R_2(x, y)$ for all $x$. Then the relation
  $R_1 \vee R_2$ is in $\cc{TFNP}$ (Megiddo--Papadimitriou).
\end{theorem}
\begin{proof}
  \leanok
  $\cc{FNP}$ is closed under disjunction of relations, and every $x$ has a
  witness for one side. With Theorem~\ref{thm:polytime-np-verifier}, every
  language in $\cc{NP} \cap \cc{coNP}$ supplies such a pair.
\end{proof}

\section{Reductions and completeness}

\begin{definition}[Polynomial-time many-one reduction]
  \label{def:reduction}
  \lean{Complexity.MapReducesPoly}
  \leanok
  \uses{def:FP, def:language}
  $L \le_p L'$ if there is $f \in \cc{FP}$ with $x \in L \iff f(x) \in L'$ for
  all $x$.
\end{definition}

\begin{lemma}[$\le_p$ is a preorder]
  \label{lem:polytime-reduction-preorder}
  \lean{Complexity.MapReducesPoly.refl, Complexity.MapReducesPoly.trans}
  \leanok
  \uses{def:reduction}
  $L \le_p L$, and $L_1 \le_p L_2 \le_p L_3$ implies $L_1 \le_p L_3$.
\end{lemma}
\begin{proof}
  \leanok
  \uses{lem:polytime-FP-comp}
  The identity and composition are in $\cc{FP}$.
\end{proof}

\begin{lemma}[Transport along reductions]
  \label{lem:polytime-reduction-transport}
  \lean{Complexity.MapReducesPoly.compl, Complexity.MapReducesPoly.mem_P, Complexity.MapReducesPoly.mem_NP, Complexity.MapReducesPoly.mem_coNP}
  \leanok
  \uses{def:reduction, def:P, def:NP, def:coNP}
  If $L_1 \le_p L_2$ then $\overline{L_1} \le_p \overline{L_2}$, and for each
  $\mathcal{C} \in \{\cc{P}, \cc{NP}, \cc{coNP}\}$, $L_2 \in \mathcal{C}$
  implies $L_1 \in \mathcal{C}$.
\end{lemma}
\begin{proof}
  \leanok
  \uses{thm:polytime-preimage}
  The same function reduces the complements, and $L_1 = f^{-1}(L_2)$.
\end{proof}

\begin{definition}[$\cc{NP}$-hard and $\cc{NP}$-complete]
  \label{def:np-complete}
  \lean{Complexity.NPHard, Complexity.NPComplete}
  \leanok
  \uses{def:reduction, def:NP}
  $L$ is $\cc{NP}$-hard if $L' \le_p L$ for every $L' \in \cc{NP}$, and
  $\cc{NP}$-complete if moreover $L \in \cc{NP}$.
\end{definition}

\begin{lemma}[Hardness transfers along reductions]
  \label{lem:polytime-np-hard-transfer}
  \lean{Complexity.NPHard.of_reduction, Complexity.NPComplete.transfer}
  \leanok
  \uses{def:np-complete}
  If $L_1$ is $\cc{NP}$-hard and $L_1 \le_p L_2$, then $L_2$ is
  $\cc{NP}$-hard; if moreover $L_2 \in \cc{NP}$, then $L_2$ is
  $\cc{NP}$-complete.
\end{lemma}
\begin{proof}
  \leanok
  \uses{lem:polytime-reduction-preorder}
  Transitivity of $\le_p$.
\end{proof}

\begin{theorem}[An $\cc{NP}$-complete language decides the questions]
  \label{thm:polytime-np-complete-collapse}
  \lean{Complexity.NPComplete.mem_P_iff_P_eq_NP, Complexity.NPComplete.mem_coNP_iff_NP_eq_coNP}
  \leanok
  \uses{def:np-complete, def:P, def:coNP}
  If $L$ is $\cc{NP}$-complete, then $L \in \cc{P} \iff \cc{P} = \cc{NP}$ and
  $L \in \cc{coNP} \iff \cc{NP} = \cc{coNP}$.
\end{theorem}
\begin{proof}
  \leanok
  \uses{lem:polytime-reduction-transport, thm:polytime-P-subset-NP}
  Every $\cc{NP}$ language pulls back into $\cc{P}$ (respectively
  $\cc{coNP}$) along its reduction to $L$.
\end{proof}

\begin{definition}[$\cc{coNP}$-hard and $\cc{coNP}$-complete]
  \label{def:polytime-conp-complete}
  \lean{Complexity.coNPHard, Complexity.coNPComplete}
  \leanok
  \uses{def:reduction, def:coNP}
  $L$ is $\cc{coNP}$-hard if $L' \le_p L$ for every $L' \in \cc{coNP}$, and
  $\cc{coNP}$-complete if moreover $L \in \cc{coNP}$. The complement of an
  $\cc{NP}$-complete language is $\cc{coNP}$-complete and conversely
  (\texttt{Complexity.NPComplete.compl}, \texttt{Complexity.coNPComplete.compl}).
\end{definition}

\section{Satisfiability and the Cook--Levin theorem}

\begin{definition}[$\mathrm{SAT}$]
  \label{def:polytime-sat}
  \lean{Complexity.SAT.language, Complexity.SAT.Witness}
  \leanok
  \uses{def:language}
  A CNF formula is a list of clauses, a clause a list of literals, and a
  literal a sign with a variable index in $\mathbb{N}$; an assignment is a bit
  string, with out-of-range variables read as false. Formulas are encoded as
  bit strings with variable indices in unary, data bits doubled, and the
  undoubled pairs $01$ and $10$ separating literals and clauses.
  $\mathrm{SAT}$ is the set of encodings of satisfiable formulas. The witness
  relation holds of $(z, \alpha)$ when $z$ encodes some $\varphi$,
  $|\alpha| \le |z| + 1$, and $\alpha$ satisfies $\varphi$.
\end{definition}

\begin{theorem}[$\mathrm{SAT} \in \cc{NP}$]
  \label{thm:polytime-sat-np}
  \lean{Complexity.SAT.pairLang_witness_mem_P, Complexity.SAT.language_mem_NP}
  \leanok
  \uses{def:polytime-sat, def:pairlang, def:NP}
  The pair language of the $\mathrm{SAT}$ witness relation is in $\cc{P}$
  (decided in quadratic time), and $\mathrm{SAT} \in \cc{NP}$.
\end{theorem}
\begin{proof}
  \leanok
  A three-work-tape deterministic verifier parses the formula and evaluates it
  at the witness; a $\mathrm{SAT}$-specialized guess-and-verify machine turns
  it into a nondeterministic decider. Unary variable indices make
  $x \in \mathrm{SAT}$ equivalent to the existence of a witness of length at
  most $|x| + 1$.
\end{proof}

\begin{theorem}[Cook--Levin]
  \label{thm:cook-levin}
  \lean{Complexity.SAT.NPComplete_language, Complexity.SAT.NPHard_language, Complexity.SAT.cookLevin_reduction}
  \leanok
  \uses{def:polytime-sat, def:np-complete, def:ntm-decides}
  $\mathrm{SAT}$ is $\cc{NP}$-complete. More precisely, if a
  nondeterministic machine with any number of work tapes decides $L$ within
  time $T$ with $T = O(n^c)$, then $L \le_p \mathrm{SAT}$.
\end{theorem}
\begin{proof}
  \leanok
  \uses{thm:polytime-sat-np}
  The machine is first converted to one with a single work tape, and $T$ is
  replaced by a dominating explicit polynomial $p$. The reduction maps $x$ to
  the tableau formula for $p(|x|)$ steps, which is satisfiable exactly when
  some path accepts. An emitter machine writes the seven encoded clause
  families of the tableau in polynomial time.
\end{proof}

\begin{corollary}[Consequences of Cook--Levin]
  \label{cor:polytime-sat-consequences}
  \lean{Complexity.SAT.coNPComplete_compl_language, Complexity.SAT.language_mem_P_iff_P_eq_NP, Complexity.SAT.language_mem_coNP_iff_NP_eq_coNP}
  \leanok
  \uses{def:polytime-sat, def:polytime-conp-complete}
  The complement of $\mathrm{SAT}$ (as a set of bit strings, so including all
  malformed strings) is $\cc{coNP}$-complete. $\mathrm{SAT} \in \cc{P} \iff
  \cc{P} = \cc{NP}$, and $\mathrm{SAT} \in \cc{coNP} \iff \cc{NP} = \cc{coNP}$.
\end{corollary}
\begin{proof}
  \leanok
  \uses{thm:cook-levin, thm:polytime-np-complete-collapse}
  Completeness dualizes under complement, and
  Theorem~\ref{thm:polytime-np-complete-collapse} applies to
  $\mathrm{SAT}$.
\end{proof}

\begin{definition}[$3\mathrm{SAT}$]
  \label{def:polytime-3sat}
  \lean{Complexity.SAT.ThreeSAT.language}
  \leanok
  \uses{def:polytime-sat}
  $3\mathrm{SAT}$ is the set of encodings, in the $\mathrm{SAT}$ codec, of
  satisfiable CNF formulas in which every clause has exactly three literals.
\end{definition}

\begin{lemma}[Tseitin splitting]
  \label{lem:polytime-tseitin}
  \lean{Complexity.SAT.ThreeSAT.reduction_mem_FP, Complexity.SAT.ThreeSAT.cnfsat_le_language}
  \leanok
  \uses{def:polytime-sat, def:polytime-3sat, def:reduction}
  There is a total function in $\cc{FP}$ witnessing
  $\mathrm{SAT} \le_p 3\mathrm{SAT}$. On the encoding of a formula it outputs
  an equisatisfiable exact-3 formula: short clauses are padded by repeating
  literals, the empty clause becomes a contradictory pair of clauses, and long
  clauses are split by a chain of fresh variables numbered above $|z|$. On a
  malformed string it outputs a fixed unsatisfiable exact-3 formula.
\end{lemma}
\begin{proof}
  \leanok
  Equisatisfiability is proved clause by clause with a threaded
  fresh-variable counter; polynomial time is a bespoke streaming transducer.
\end{proof}

\begin{theorem}[$3\mathrm{SAT}$ is $\cc{NP}$-complete]
  \label{thm:polytime-3sat-complete}
  \lean{Complexity.SAT.ThreeSAT.NPComplete_language, Complexity.SAT.ThreeSAT.language_mem_NP}
  \leanok
  \uses{def:polytime-3sat, def:np-complete}
  $3\mathrm{SAT}$ is $\cc{NP}$-complete.
\end{theorem}
\begin{proof}
  \leanok
  \uses{thm:polytime-sat-np, lem:polytime-tseitin, thm:cook-levin, lem:polytime-np-hard-transfer}
  Membership: a witness is a $\mathrm{SAT}$ witness whose formula also passes
  the exact-3 syntax check, a regular language decided by a finite-state
  machine. Hardness: compose Cook--Levin with Tseitin splitting.
\end{proof}

\begin{theorem}[Padded circuit satisfiability is in $\cc{NP}$]
  \label{thm:polytime-circuit-sat-np}
  \lean{Complexity.CircuitSAT.language, Complexity.CircuitSAT.language_mem_NP, Complexity.CircuitSAT.extensionLanguage_mem_NP}
  \leanok
  \uses{def:NP, def:circuit, def:pair}
  A query is a pair $\langle c, r \rangle$ of a tagged circuit-family code
  $c$ and a ruler $r$ whose length fixes the assignment width; it is accepted
  when some $w$ with $|w| = |r|$ makes the serialized circuit evaluator return
  true on $c$ and $w$. This language is in $\cc{NP}$, and so is the extension
  variant with queries $\langle c, \langle u, r \rangle \rangle$, accepted when
  some $w$ with $|w| = |r|$ makes $c$ accept $u\,w$.
\end{theorem}
\begin{proof}
  \leanok
  \uses{lem:polytime-guess-verify}
  The verified serialized circuit evaluator decides the paired witness
  language in polynomial time, and witnesses are linearly bounded by the query.
  $\cc{NP}$-hardness of this language is not yet stated.
\end{proof}

\section{More NP-complete problems}

\begin{theorem}[Formula satisfiability]
  \label{thm:polytime-formula-sat}
  \uses{def:polytime-sat, def:np-complete, def:reduction}
  Fix a concrete prefix codec for Boolean formulas over $\wedge$, $\vee$,
  $\neg$ and variables, and let $\mathrm{FSAT}$ be the set of encodings of
  satisfiable formulas, malformed strings excluded. Then
  $\mathrm{SAT} \le_p \mathrm{FSAT}$ (a CNF is a formula) and
  $\mathrm{FSAT} \le_p \mathrm{SAT}$ (Tseitin encoding with one fresh variable
  per gate), so $\mathrm{FSAT}$ is $\cc{NP}$-complete.
\end{theorem}
\begin{proof}
  \uses{thm:cook-levin, lem:polytime-np-hard-transfer, lem:polytime-guess-verify, lem:polytime-fp-combinators}
  Membership by guess and verify with a formula evaluator; both reductions are
  streaming transducers, most easily written through
  Lemma~\ref{lem:polytime-fp-combinators}.
\end{proof}

\begin{definition}[Graph problems]
  \label{def:polytime-graph-problems}
  \uses{def:pair, def:language}
  Fix a codec for pairs $\langle G, k \rangle$ of a finite simple graph
  (vertex count in unary and a row-major adjacency matrix) and a threshold $k$
  in unary, with malformed strings as no-instances.
  $\mathrm{INDSET}$, $\mathrm{CLIQUE}$, and $\mathrm{VC}$ are the sets of
  $\langle G, k \rangle$ such that $G$ has an independent set of size at least
  $k$, a clique of size at least $k$, and a vertex cover of size at most $k$,
  respectively.
\end{definition}

\begin{theorem}[$\mathrm{INDSET}$ is $\cc{NP}$-complete]
  \label{thm:polytime-indset-complete}
  \uses{def:polytime-graph-problems, def:np-complete}
  $\mathrm{INDSET}$ is $\cc{NP}$-complete.
\end{theorem}
\begin{proof}
  \uses{thm:polytime-3sat-complete, lem:polytime-guess-verify, lem:polytime-np-hard-transfer, lem:polytime-fp-combinators}
  Membership by guess and verify on a vertex subset. Hardness by
  $3\mathrm{SAT} \le_p \mathrm{INDSET}$: one triangle per clause, edges between
  complementary literals, and $k$ the number of clauses.
\end{proof}

\begin{theorem}[$\mathrm{CLIQUE}$ and $\mathrm{VC}$ are $\cc{NP}$-complete]
  \label{thm:polytime-clique-vc-complete}
  \uses{def:polytime-graph-problems, def:np-complete}
  $\mathrm{CLIQUE}$ and $\mathrm{VC}$ are $\cc{NP}$-complete.
\end{theorem}
\begin{proof}
  \uses{thm:polytime-indset-complete, lem:polytime-np-hard-transfer, lem:polytime-guess-verify}
  $\mathrm{INDSET} \le_p \mathrm{CLIQUE}$ by complementing the edge set, and
  $\mathrm{INDSET} \le_p \mathrm{VC}$ by $\langle G, k \rangle \mapsto
  \langle G, n - k \rangle$, since $S$ is independent exactly when its
  complement is a cover.
\end{proof}

\section{Open directions}

Two refactors do not fit as nodes.

\begin{itemize}
  \item \textbf{Re-proving bespoke machines through Cobham.} The
    $\mathrm{SAT}$ verifier and guess-and-verify machine behind
    Theorem~\ref{thm:polytime-sat-np} take roughly 11--12 thousand lines, the
    Cook--Levin emitter about 6.7 thousand, and the Tseitin transducer of
    Lemma~\ref{lem:polytime-tseitin} about 6 thousand. With the combinators
    of Lemma~\ref{lem:polytime-fp-combinators}, each should become a short
    Cobham program plus a correctness proof about ordinary Lean functions, as
    already done for the Sipser--Lautemann matrix.
  \item \textbf{Surface placement.} Lemma~\ref{lem:polytime-guess-verify}
    (\texttt{mem\_NP\_of\_poly\_witness}) lives in
    \texttt{Classes/PCP/Internal} and Lemma~\ref{lem:polytime-iterate} in
    \texttt{Classes/P/Cobham/Internal}; both are general tools and belong in
    surface modules under \texttt{Classes/NP} and \texttt{Classes/P}.
\end{itemize}
